\chapter{Step 4: the submission}
\label{ch:submission}
\whereami{5}

\begin{tldr}
\begin{itemize}
\item The submission notebook runs \textbf{offline} on hidden data: everything it needs is attached as datasets.
\item \textbf{Safety net first}: write a valid answer file before doing anything risky, then improve it.
\item Each test study is prepared by the \textbf{same} \code{knee\_preproc.py}, predicted by every attached model
      group (with mirror TTA), and the groups are \textbf{rank-blended}.
\item To use a public notebook, we wrote a ``member cell'' that blends our ranks into its predictions.
\end{itemize}
\end{tldr}

\section{The flow}

\begin{figure}[H]
\centering
\begin{tikzpicture}[x=0.9mm,y=1mm,
  st/.style={draw=accent, fill=keyfill, rounded corners=1.5mm, minimum height=12mm, align=center,
             font=\scriptsize, text width=24mm},
  ar/.style={-{Stealth[length=1.8mm]}, line width=0.7pt, accent}]
  \node[st, draw=wordline, fill=wordfill] (a) at (0,0) {\textbf{1. safety net}\\ write a valid\\ {\ttfamily\tiny submission.csv}};
  \node[st] (b) at (30,0) {\textbf{2. offline installs}\\ JPEG decoders\\ from attached wheels};
  \node[st] (c) at (60,0) {\textbf{3. load}\\ {\ttfamily\tiny knee\_preproc.py}\\ + model groups};
  \node[st] (d) at (95,8) {\textbf{4a. prepare}\\ 32 studies\\ (4 CPU workers)};
  \node[st] (e) at (95,-12) {\textbf{4b. predict}\\ every group,\\ normal + mirrored};
  \node[st, draw=wordline, fill=wordfill] (f) at (130,-2) {\textbf{5. rewrite}\\ {\ttfamily\tiny submission.csv}\\ every few chunks};
  \draw[ar] (a) -- (b); \draw[ar] (b) -- (c); \draw[ar] (c) -- (d.west);
  \draw[ar] (d) -- (e); \draw[ar] (e.east) -- (f);
  \draw[ar, dashed] (f.north) |- (d.east) node[pos=0.75, above, font=\scriptsize] {next chunk};
\end{tikzpicture}
\caption{The submission notebook. Steps 4a--5 repeat until all \ca 1,300 test studies are done or the time limit
approaches.}
\end{figure}

\section{Safety net first}

\begin{keyidea}
A code competition punishes crashes harshly: a crash means no score, discovered hours later. So the first thing the
notebook does is write a complete, valid answer file (one guess per finding, the training prevalence). Every study is
then processed inside its own \code{try/except}; a study that fails keeps its guess, and the run goes on.
\end{keyidea}

\begin{watchout}
Two of our runs of the public-notebook version failed only because a stray one-line cell (containing the text
\code{[code]}, a paste accident) sat before our cell. A failed commit run means a failed scoring run, but you only
learn that \ca 6 hours later. \textbf{Check the cell list before you submit.}
\end{watchout}

\section{Offline installs}

Some sites store their DICOM pixels compressed (JPEG Lossless, JPEG 2000). Decoding them needs extra Python packages,
and the scoring run has no internet. We downloaded the packages' installation files (\term{wheels}) once into a Kaggle
dataset, \code{knee-dicom-wheels}, and the notebook installs them from there.

\section{Model groups and rank blending}

Every attached model dataset that contains a \code{train\_config.json} is one \term{group} (for example ``v5 B0'',
``v5 ConvNeXt-Nano''). Within a group, the 5 fold models and the 2 views (normal + mirrored) are averaged. Between
groups, predictions are combined \textbf{by rank}, per finding:

\begin{figure}[H]
\centering
\begin{tikzpicture}[x=1mm,y=1mm]
  \matrix[matrix of nodes, nodes={font=\footnotesize, minimum width=17mm, minimum height=6mm, anchor=center},
          column 1/.style={nodes={font=\footnotesize\bfseries}}, row 1/.style={nodes={font=\scriptsize\bfseries, text=muted}}] (m) {
    study & model 1 & model 2 & \textcolor{muted}{prob.\ avg.} & rank 1 & rank 2 & \textbf{rank avg.} \\
    A & 0.95 & 0.52 & 0.735 & 4 & 3 & \textbf{3.5} \\
    B & 0.10 & 0.55 & 0.325 & 2 & 4 & \textbf{3.0} \\
    C & 0.40 & 0.50 & 0.450 & 3 & 2 & \textbf{2.5} \\
    D & 0.05 & 0.49 & 0.270 & 1 & 1 & \textbf{1.0} \\
  };
  \draw[grid] (m.north west) ++(0,-7) -- ++(119,0);
  \node[font=\scriptsize, text=warnline, align=left, anchor=west] at (m.south west) [yshift=-6mm]
    {probability average: order A, C, B, D = model 1's order; the cautious model 2 is ignored};
  \node[font=\scriptsize, text=wordline, align=left, anchor=west] at (m.south west) [yshift=-11mm]
    {rank average: order A, B, C, D; both models count equally};
\end{tikzpicture}
\caption{Why rank blending. Model 2 is cautious (all its predictions are close to 0.5), so averaging probabilities lets
model 1 decide alone. Ranks put both on the same scale. Since AUC only cares about order, nothing is lost.}
\label{fig:rankblend}
\end{figure}

\code{BLEND\_WEIGHTS} can give a group a different weight per finding (for example, weight 0 for the older v3 group on
ACL, MCL and PF OA, where its labels were known to be bad).

\section{Plugging into a public notebook: the ``member cell''}

By early October the public leaderboard was crowded at \ca 0.94 by copies (\term{forks}) of one public notebook that
ran several strong published models (Chapter~\ref{ch:experiments}). We forked it and added one cell of our own:

\begin{figure}[H]
\centering
\begin{tikzpicture}[x=1mm,y=1mm,
  cell/.style={draw=grid, fill=white, rounded corners=1mm, minimum width=36mm, minimum height=8mm, font=\scriptsize,
               align=center},
  ours/.style={cell, draw=us, fill=keyfill, line width=0.9pt}]
  \node[cell] (c1) at (0,24) {public cells: their models};
  \node[cell] (c2) at (0,12) {staged predictions (ranks)};
  \node[ours, minimum height=14mm, text width=34mm] (c3) at (0,-3) {\textbf{our member cell}\\ back up $\cdot$ run our\\ models $\cdot$ blend ranks};
  \node[cell] (c4) at (0,-19) {publish \code{submission.csv}};
  \draw[-{Stealth}, muted] (c1) -- (c2); \draw[-{Stealth}, muted] (c2) -- (c3); \draw[-{Stealth}, muted] (c3) -- (c4);
  \node[font=\scriptsize, align=left, anchor=west] at (24,6) {new rank $= (1-w)\cdot$ public rank $+\ w\cdot$ our rank\\[1mm]
     $w$ = \code{OUR\_WEIGHT}, we tried 0, 0.25 and 0.4};
  \node[font=\scriptsize, align=left, anchor=west, text=warnline] at (24,-6) {any error or time-out: restore the backup,\\
     so the public result is never lost};
\end{tikzpicture}
\caption{The member cell sits just before the public notebook's final ``publish'' cell. It only touches the staged
prediction file and restores it if anything goes wrong.}
\label{fig:member}
\end{figure}

\begin{innumbers}
\begin{itemize}
\item A commit run of our own notebook (3 sample studies) took about 1--8 minutes, mostly loading models.
\item Scoring a submission of the public notebook took \ca 6 hours before the result (or an error) appeared.
\item Our member cell had a time limit (1.5--2 hours, or 20 minutes before the public notebook's own budget).
\end{itemize}
\end{innumbers}

\begin{checkpoint}
\begin{itemize}
\done I can explain the safety-net pattern and why every study gets its own \code{try/except}.
\done I can rank-blend two models by hand.
\quiz Q10.1: Why must the test images be prepared with exactly the training cache's code?
\quiz Q10.2: Two models disagree completely on one finding. What does the rank average look like?
\quiz Q10.3: What does the member cell do if our models crash halfway?
\end{itemize}
\end{checkpoint}
